\section{问题重述}

超声波流量计利用斜穿管道的声道中超声波顺流和逆流传播时间的差异来估计流体流速。探头固定在管壁上，声波相对管壁的速度需叠加水流拖曳作用——顺流时声波被水流推动，传播时间缩短；逆流时声波被水流阻挡，传播时间延长。通过精确测量这一时间差，可计算声道方向的水流速度分量。

单声道流量计仅能测量该声道路径上的线平均速度，而管道截面上的流速分布并不均匀——管中心流速高、靠近管壁处流速低。这种截面速度分布的非均匀性使得单声道测量结果无法代表整个截面的面平均流速。工业级超声波流量计通过在截面不同高度布设多条声道，对各声道测量值进行加权求和，以数值积分的方式逼近面平均流量。声道权重的选择直接影响积分精度——权重需反映各声道所代表的环带截面积占总面积的比例，这是多声道流量计设计的核心问题。

当管道上游存在弯头、阀门、缩径等扰流件时，流速剖面发生畸变——可能出现非对称分布、旋流、偏流等现象。这些剖面畸变使固定权重积分产生系统性偏差，导致测量精度恶化。本文所涉及的五声道超声波流量计在八种不同扰流状态下运行，每次测量以窗口为单位，每窗口包含五条声道的等效观测量及若干反映流速剖面形状和窗口稳定性的在线特征。

本文需解决以下四个问题：

\textbf{问题1（基础模型验证）}：从顺逆流传播时间模型出发，推导单声道流速公式，说明多声道积分的必要性，给出一种五声道加权积分方法并解释权重物理含义，对比附件中四套积分方案的误差表现。

\textbf{问题2（无扰流建模）}：仅使用无扰流数据（10个窗口），建立多声道流量估计模型，要求说明输入变量、参数估计方法和防过拟合措施。

\textbf{问题3（扰流识别与补偿）}：分析各在线特征对扰流的判别能力，对八种扰流状态进行聚类归并，构造不依赖扰流编号的在线识别规则，确定是否需要分流量点进行扰流补偿。

\textbf{问题4（最终达标模型）}：综合前三问成果建立满足工程要求的最终模型，报告各项评价指标并分析是否达标。
